\begin{ex}
Trong phép chia $17 : 3$, số dư là bao nhiêu?
\shortans{2}
\loigiai{
$17 = 3 \times 5 + 2$, dư 2.
}
\end{ex}

\begin{ex}
Trong phép chia $23 : 4$, số dư là bao nhiêu?
\shortans{3}
\loigiai{
$23 = 4 \times 5 + 3$, dư 3.
}
\end{ex}

\begin{ex}
Trong phép chia $29 : 5$, thương là bao nhiêu?
\shortans{5}
\loigiai{
$29 = 5 \times 5 + 4$, thương là 5.
}
\end{ex}

\begin{ex}
Trong phép chia $38 : 6$, thương là bao nhiêu?
\shortans{6}
\loigiai{
$38 = 6 \times 6 + 2$, thương là 6.
}
\end{ex}

\begin{ex}
Trong phép chia $45 : 7$, số dư là bao nhiêu?
\shortans{3}
\loigiai{
$45 = 7 \times 6 + 3$, dư 3.
}
\end{ex}

\begin{ex}
Trong phép chia $53 : 8$, số dư là bao nhiêu?
\shortans{5}
\loigiai{
$53 = 8 \times 6 + 5$, dư 5.
}
\end{ex}

\begin{ex}
Trong phép chia $67 : 9$, số dư là bao nhiêu?
\shortans{4}
\loigiai{
$67 = 9 \times 7 + 4$, dư 4.
}
\end{ex}

\begin{ex}
Trong một phép chia có số chia là 6, số dư lớn nhất có thể có là bao nhiêu?
\shortans{5}
\loigiai{
Số dư luôn nhỏ hơn số chia, lớn nhất là $6 - 1 = 5$.
}
\end{ex}

\begin{ex}
Trong một phép chia có số chia là 8, số dư lớn nhất có thể có là bao nhiêu?
\shortans{7}
\loigiai{
$8 - 1 = 7$.
}
\end{ex}

\begin{ex}
Trong một phép chia cho 5, số dư có thể nhận bao nhiêu giá trị khác nhau (kể cả số dư bằng 0)?
\shortans{5}
\loigiai{
Các số dư có thể là 0, 1, 2, 3, 4 (5 giá trị).
}
\end{ex}

\begin{ex}
Có 19 bông hoa cắm đều vào các lọ, mỗi lọ 4 bông hoa. Hỏi cắm được nhiều nhất bao nhiêu lọ hoa?
\shortans{4}
\loigiai{
$19 : 4 = 4$ dư 3 lọ.
}
\end{ex}

\begin{ex}
Có 25 cái kẹo chia đều cho 4 bạn. Hỏi mỗi bạn được nhiều nhất bao nhiêu cái kẹo?
\shortans{6}
\loigiai{
$25 : 4 = 6$ dư 1.
}
\end{ex}

\begin{ex}
Có 31 học sinh đi đò, mỗi chuyến chở được 6 học sinh. Hỏi cần ít nhất bao nhiêu chuyến đò để chở hết số học sinh?
\shortans{6}
\loigiai{
$31 : 6 = 5$ dư 1, cần ít nhất $5 + 1 = 6$ chuyến.
}
\end{ex}

\begin{ex}
Một lớp có 34 học sinh. Mỗi bàn ngồi 2 học sinh. Hỏi cần ít nhất bao nhiêu bàn cho cả lớp ngồi?
\shortans{17}
\loigiai{
$34 : 2 = 17$ bàn (chia hết).
}
\end{ex}

\begin{ex}
Một phép chia có số chia là 7, thương là 8 và số dư là 5. Số bị chia là bao nhiêu?
\shortans{61}
\loigiai{
$8 \times 7 + 5 = 56 + 5 = 61$.
}
\end{ex}

\begin{ex}
Một phép chia có số chia là 9, thương là 6 và số dư lớn nhất có thể. Số bị chia là bao nhiêu?
\shortans{62}
\loigiai{
Số dư lớn nhất là 8: Số bị chia = $6 \times 9 + 8 = 54 + 8 = 62$.
}
\end{ex}

\begin{ex}
Một phép chia có số chia là 5, thương là 12 và số dư là 3. Số bị chia là bao nhiêu?
\shortans{63}
\loigiai{
$12 \times 5 + 3 = 60 + 3 = 63$.
}
\end{ex}

\begin{ex}
Một cuộn dây dài 43 m được cắt thành các đoạn dài 5 m. Hỏi cắt được bao nhiêu đoạn và thừa ra mấy mét? (Điền số mét thừa).
\shortans{3}
\loigiai{
$43 : 5 = 8$ dư 3 m.
}
\end{ex}

\begin{ex}
Có 50 chiếc bút xếp vào các hộp, mỗi hộp đựng 8 chiếc. Hỏi còn thừa mấy chiếc bút?
\shortans{2}
\loigiai{
$50 : 8 = 6$ dư 2 chiếc.
}
\end{ex}

\begin{ex}
Tính thương trong phép chia $49 : 6$.
\shortans{8}
\loigiai{
$49 = 6 \times 8 + 1$, thương là 8.
}
\end{ex}

\begin{ex}
Số nào chia cho 4 được thương là 9 và dư 2?
\shortans{38}
\loigiai{
$9 \times 4 + 2 = 38$.
}
\end{ex}

\begin{ex}
Số nào chia cho 7 được thương là 5 và dư 6?
\shortans{41}
\loigiai{
$5 \times 7 + 6 = 41$.
}
\end{ex}

\begin{ex}
Trong phép chia cho 9, có bao nhiêu số dư khác 0 có thể có?
\shortans{8}
\loigiai{
Các số dư từ 1 đến 8 (8 giá trị).
}
\end{ex}

\begin{ex}
May mỗi bộ quần áo cần 3 m vải. Có 26 m vải thì may được nhiều nhất bao nhiêu bộ quần áo?
\shortans{8}
\loigiai{
$26 : 3 = 8$ dư 2 m, may được 8 bộ.
}
\end{ex}

\begin{ex}
Số dư của phép chia $75 : 8$ là bao nhiêu?
\shortans{3}
\loigiai{
$75 = 8 \times 9 + 3$, số dư là 3.
}
\end{ex}
